% !TEX program = lualatex
\documentclass[11pt,a4paper]{article}
\usepackage{luatexja}
\usepackage{amsmath,amssymb,amsthm,amsfonts}
\usepackage{geometry}
\geometry{margin=1in}

% 定理環境の設定
\theoremstyle{definition}
\newtheorem{definition}{定義}[section]
\newtheorem{theorem}{定理}[section]
\newtheorem{lemma}{補題}[section]
\newtheorem{proof_topological}{証明の位相幾何学的構成}

\title{真理保存性の不完全性定理の位相幾何学的モデル：\\メビウス多様体における自己言及構造の解消}
\author{千葉将臣}
\date{2026-04-30}

\begin{document}

\maketitle

\begin{abstract}
本稿は、形式体系における「真理保存性」の要請とそれに伴う限界定理的現象（不完全性）を、一次元多様体上のファイバー束の位相的性質を用いてモデル化する。厳密な自己同一性（$A=A$）を要請する真理保存体系を直積空間（円柱）として、同型性（$A \cong A$）を許容する体系を非向き付け可能空間（メビウスの帯）として構成する。これにより、古典的体系における自己言及の決定不能性が、大域的な位相的障害（非連結性）として幾何学的に解釈可能であることを示す。なお、本稿はGödelの不完全性定理の構文論的な一般化を意図するものではなく、限界定理的現象の位相的モデルの構成例を提示するものである。
\end{abstract}

\section{導入と幾何学的モデルの定義}

形式的推論体系における「真理値」の推移を、基空間 $S^1$（推論の反復プロセス）上のファイバー束としてモデル化する。状態空間のファイバーを $F = \{-1, 1\}$ とし、$1$ を真（$A$）、$-1$ を偽（$\neg A$）に対応させる。

\begin{definition}[体系 $T_{=}$：厳密な真理保存性]
厳密な自己同一性 $A=A$ を維持する体系 $T_{=}$ の状態空間 $E_{=}$ を、自明なファイバー束（円柱の境界）として定義する。
$$ E_{=} = S^1 \times \{-1, 1\} $$
この空間は位相的に2つの交わらない閉曲線 $S^1 \sqcup S^1$ と同相であり、非連結である（$\pi_0(E_{=}) = 2$）。
\end{definition}

\begin{definition}[体系 $T_{\cong}$：同型性による真理保存]
全体としての同型性 $A \cong A$ を維持する体系 $T_{\cong}$ の状態空間 $E_{\cong}$ を、非向き付け可能なファイバー束（メビウスの帯の境界）として定義する。すなわち、$[0, 1] \times \{-1, 1\}$ に対して同値関係 $(0, y) \sim (1, -y)$ を入れた商空間とする。
$$ E_{\cong} = \frac{[0, 1] \times \{-1, 1\}}{(0, y) \sim (1, -y)} $$
この空間は位相的に1つの閉曲線 $S^1$ と同相であり、連結である（$\pi_0(E_{\cong}) = 1$）。
\end{definition}

\section{限界定理的現象のトポロジカルな定式化}

体系内部における自己言及的否定命題 $G$（「本プロセスは自らの否定状態へ到達する」）の検証を、基空間に沿った連続的な推論経路 $\gamma(t)$ として定式化する。

\begin{definition}[命題 $G$ の連続写像による表現]
初期状態 $\gamma(0) = (0, 1)$ から出発し、基空間を一周した結果として自身の否定状態 $\gamma(1) = (1, -1)$ に到達する連続写像 $\gamma: [0, 1] \to E$ が存在するとき、命題 $G$ はその空間上で構成可能であるとする。
\end{definition}

\begin{theorem}[真理保存の境界と位相的障害]
\label{thm:mobius_incompleteness}
(1) 厳密な自己同一性を要請する体系 $T_{=}$ において、自己言及経路 $G$ は構成不能（到達不能）である。\\
(2) 同型性を要請する体系 $T_{\cong}$ において、経路 $G$ は構成可能であり、その構造は空間の基本群の生成元に一致する。
\end{theorem}

\begin{proof_topological}
(1) 体系 $T_{=}$ （円柱）における位相的障害：\\
空間 $E_{=} = S^1 \times \{-1, 1\}$ において、初期状態 $y=1$ を持つ連結成分と、状態 $y=-1$ を持つ連結成分は完全に分離している。したがって、$\gamma(0) = (0, 1)$ を満たす連続写像は $\gamma(t) = (t, 1)$ と定値に保たれざるを得ず、$\gamma(1) = (1, -1)$ を満たす経路は存在しない。これは、自己同一性という位相的境界（非連結性）により、否定への到達が構造的に阻まれることを示している。

(2) 体系 $T_{\cong}$ （メビウス多様体）における構造的解消：\\
空間 $E_{\cong}$ は同値関係 $(0, 1) \sim (1, -1)$ を持つ。経路 $\gamma(t) = (t, 1)$ は局所座標系において「真（$y=1$）」を維持して進むが、基空間を一周し $t \to 1$ の極限をとると、空間全体の同値関係により、到達点は幾何学的に $(1, -1)$ と同一視される。
すなわち、局所的な「真」の連続的推移が、大域的には「自らの否定状態への到達」として実現される。この閉曲線 $\gamma$ は、空間 $E_{\cong}$ の基本群 $\pi_1(E_{\cong}) \cong \mathbb{Z}$ の生成元をなす。
\end{proof_topological}

\section{結論}

古典的論理体系における自己言及のパラドックスや決定不能性といった限界定理的現象は、本モデルによれば、真理保存性を「空間の非連結性（自明な直積空間）」として定義したことに伴う位相的制約として幾何学的に解釈される。真理保存の条件を連続的な同型性（非向き付け可能多様体）へと拡張することにより、到達不能であった命題 $G$ は、空間が持つ大域的な位相的性質（ねじれを伴う連結性）を記述する構造へと変化する。本モデルは、形式体系の限界と次元的拡張（真理保存性の緩和）の関係性を、トポロジーの視点から可視化する直観的基盤を提供するものである。
\end{document}
